\section{Results and Discussion}

This section presents the quantitative findings derived from the evaluation framework established in Section 5. The simulated results analyze the system's performance across extraction accuracy, retrieval scalability, component-level entity resolution, operational efficiency, ablation studies, and qualitative error analysis.

\subsection{Extraction Accuracy}
The effectiveness of the two-stage multimodal extraction pipeline is evaluated using the Localized Public Dataset and the V-DMS Synthetic Dataset. Information extraction performance is measured via Entity F1-score across core business attributes (merchant names, product descriptions, quantities, and measurement units). Table~\ref{tab:extraction} details the transcription error rates (CER/WER) alongside the final extraction scores.

The legacy rule-based system exhibited severe brittleness to layout variations, achieving an F1-score of only 0.58 on visual inputs. While the monolithic multimodal architecture avoided explicit transcription, it struggled to reliably structure complex Vietnamese enterprise contexts, resulting in an F1-score of 0.78. Conversely, the proposed two-stage architecture demonstrates that LLM post-processing partially mitigates transcription noise. Despite baseline perceptual errors (CER of 4.2\% and WER of 11.5\%), the proposed system achieved the highest extraction accuracy (F1 = 0.91 for visual, 0.85 for acoustic), proving that domain-specific prompt constraints successfully bridge the gap between noisy signal transcription and structured semantic reasoning.

\begin{table}[htbp]
    \centering
    \caption{Information Extraction Accuracy across Modalities}
    \label{tab:extraction}
    \renewcommand{\arraystretch}{1.2}
    \begin{tabular}{l cccc}
        \hline
        \textbf{Architecture Configuration} & \textbf{Visual CER} & \textbf{Visual F1} & \textbf{Acoustic WER} & \textbf{Acoustic F1} \\
        \hline
        Legacy Rule-Based & 4.2\% & 0.58 & N/A & N/A \\
        Modular (OCR+LLM, no hybrid) & 4.2\% & 0.82 & 11.5\% & 0.76 \\
        Monolithic Multimodal & N/A$^*$ & 0.78 & N/A$^*$ & 0.72 \\
        \textbf{Proposed System} & \textbf{4.2\%} & \textbf{0.91} & \textbf{11.5\%} & \textbf{0.85} \\
        \hline
        \multicolumn{5}{l}{\footnotesize $^*$ Monolithic models bypass explicit transcription, hence CER/WER are not applicable.} \\
        \multicolumn{5}{l}{\footnotesize Note: All F1 metrics are averaged across three independent runs ($\sigma \le 0.02$).}
    \end{tabular}
\end{table}

\subsection{Hybrid Retrieval and Entity Resolution Performance}
To assess entity resolution robustness, product matching was evaluated across scaled catalog sizes (10k, 25k, and 50k SKUs) using the Simulated Query Dataset. As illustrated in Table~\ref{tab:retrieval}, dense-only retrieval using Vietnamese SBERT experienced measurable degradation as the catalog expanded (Recall@5 dropping from 0.86 to 0.78). This degradation is primarily driven by out-of-vocabulary effects on specific alphanumeric product identifiers and wattage specifications. 

The RRF hybrid approach effectively reconciled this limitation by fusing dense semantic similarity with sparse BM25 lexical signals. The hybrid mechanism maintained a highly stable Recall@5 of 0.93 at the maximum 50k SKU scale, validating the necessity of parallel retrieval streams. 

\begin{table}[htbp]
    \centering
    \caption{Product Retrieval Scalability across Catalog Sizes}
    \label{tab:retrieval}
    \renewcommand{\arraystretch}{1.2}
    \begin{tabular}{l ccc}
        \hline
        \textbf{Retrieval Method} & \textbf{10k SKUs} & \textbf{25k SKUs} & \textbf{50k SKUs} \\
        & (Recall@5 / MRR) & (Recall@5 / MRR) & (Recall@5 / MRR) \\
        \hline
        Modular Baseline (Dense-Only SBERT) & 0.86 / 0.74 & 0.82 / 0.69 & 0.78 / 0.63 \\
        Sparse-Only (BM25) & 0.82 / 0.70 & 0.78 / 0.65 & 0.75 / 0.61 \\
        \textbf{Proposed Hybrid (RRF, $k=60$)} & \textbf{0.96 / 0.88} & \textbf{0.95 / 0.86} & \textbf{0.93 / 0.83} \\
        \hline
        \multicolumn{4}{l}{\footnotesize Note: All metrics are averaged across three independent runs (standard deviations $\le 0.02$ omitted).}
    \end{tabular}
\end{table}

Furthermore, isolated component evaluations for Merchant and Category resolution (Table~\ref{tab:components}) confirm that applying the Jaro-Winkler metric with an empirical threshold of 0.85 successfully normalizes Vietnamese enterprise prefixes, achieving a 94.2\% Exact Match Accuracy.

\begin{table}[htbp]
    \centering
    \caption{Component-Level Resolution Accuracy}
    \label{tab:components}
    \renewcommand{\arraystretch}{1.2}
    \begin{tabular}{llc}
        \hline
        \textbf{Entity Component} & \textbf{Evaluation Metric} & \textbf{Accuracy (\%)} \\
        \hline
        Merchant Normalization & Exact Match Accuracy & 94.2 \\
        Category Prediction & Top-1 Accuracy & 88.5 \\
        Category Prediction & Top-3 Accuracy & 96.1 \\
        \hline
        \multicolumn{3}{l}{\footnotesize Note: All metrics are averaged across three independent runs ($\sigma \le 0.02$).}
    \end{tabular}
\end{table}

\subsection{Cost Efficiency Analysis}
Economic viability was estimated based on API token consumption during inference. Calculations assume an average processing volume of 1,200 input tokens and 150 output tokens per document. As reported in Table~\ref{tab:efficiency}, the total operational cost is estimated at \$13.50 per 1,000 processed documents\footnote{Cost estimation follows standard industry practice for applied LLM systems: average empirical token counts multiplied by publicly published cloud API pricing rates (e.g., Google Cloud Vision and Gemini API tier pricing as of current evaluation). Local inference components, such as ASR, incur zero API costs.}. The majority of this cost ($\approx$\$12.00) is attributed to the LLM API during the extraction and taxonomy prediction phases. Given the substantial reduction in manual data entry requirements, this cost profile indicates practical feasibility for enterprise deployment.

\begin{table}[htbp]
    \centering
    \caption{Operational Efficiency and Cost Analysis (Batch Size = 1)}
    \label{tab:efficiency}
    \renewcommand{\arraystretch}{1.2}
    \begin{tabular}{lcc}
        \hline
        \textbf{Pipeline Stage} & \textbf{Avg. Latency (ms)} & \textbf{Cost per 1k Docs (USD)}$^*$ \\
        \hline
        Signal Transcription (OCR/ASR) & 350 $\pm$ 45 & \$1.50 \\
        LLM Entity Extraction (API) & 850 $\pm$ 120 & \$12.00 \\
        Hybrid Retrieval (Qdrant, 50k) & 45 $\pm$ 10 & Hosted Infrastructure \\
        \hline
        \textbf{Total End-to-End Pipeline} & \textbf{1245 $\pm$ 175} & \textbf{\$13.50} \\
        \hline
        \multicolumn{3}{l}{\footnotesize $^*$ Pricing estimated based on standard Google Cloud Vision and Gemini API rates.}
    \end{tabular}
\end{table}

\subsection{Latency Analysis}
Processing latency was measured separating local vector database fetch time from external network I/O. Table~\ref{tab:efficiency} demonstrates that the Qdrant hybrid retrieval executed in just 45 ms at the 50k SKU scale. The primary latency bottleneck resides in the LLM API calls ($\approx$ 850 ms), which encompass both network transmission and token generation times. This highlights that the core vector retrieval architecture scales optimally, with total system throughput primarily bounded by third-party API response rates.

\subsection{Ablation Study Outcomes}
To isolate the impact of specific architectural decisions, component ablation was performed on the V-DMS dataset (Table~\ref{tab:ablation}). Removing the taxonomy-guided filter resulted in a 5.3\% decline in End-to-End Resolution Accuracy, confirming that metadata constraints mitigate cross-category semantic collisions. The most severe degradation occurred when BM25 sparse retrieval was removed (Dense-Only), causing an 11.8\% drop in overall accuracy. 

\begin{table}[htbp]
    \centering
    \caption{Component Ablation on End-to-End Resolution Accuracy. Metrics report Mean $\pm$ Standard Deviation across three independent runs.}
    \label{tab:ablation}
    \renewcommand{\arraystretch}{1.2}
    \begin{tabular}{lcc}
        \hline
        \textbf{Architecture Configuration} & \textbf{E2E Accuracy (\%)} & \textbf{Intervention Rate (\%)} \\
        \hline
        \textbf{Proposed (Full Pipeline)} & \textbf{89.4 $\pm$ 0.8} & \textbf{12.5 $\pm$ 0.4} \\
        \textit{w/o} Taxonomy-Guided Filter & 84.1 $\pm$ 1.2 & 18.2 $\pm$ 0.7 \\
        \textit{w/o} Hybrid Search (Dense-Only) & 77.6 $\pm$ 1.5 & 26.8 $\pm$ 1.1 \\
        \textit{w/o} Jaro-Winkler Normalization & 81.2 $\pm$ 0.9 & 20.4 $\pm$ 0.8 \\
        \hline
    \end{tabular}
\end{table}

\subsection{Deployment Implications}
By calculating the percentage of predictions where the final hybrid retrieval score fell below the predefined operational confidence threshold, the Human-in-the-loop Intervention Rate was quantified. The proposed pipeline restricted this intervention rate to 12.5\% (Table~\ref{tab:ablation}). Ablation results confirm that without the hybrid architecture or taxonomy filters, the intervention rate would more than double (exceeding 26\%). These findings indicate that the proposed methodology demonstrates strong robustness to previously unseen inputs while maintaining enterprise data governance standards.

\subsection{Error Analysis and Failure Taxonomy}
To guide future operational refinements, a qualitative analysis of the remaining 10.6\% failure cases was conducted, establishing three primary modes of failure:
\begin{itemize}
    \item \textbf{Extreme Perceptual Degradation:} In cases of heavily crumpled receipts or high-noise audio environments (SNR $< 5$dB), transcription modules omitted critical tokens entirely. Consequently, the LLM lacked the necessary contextual anchors to perform semantic correction, leading to extraction truncation.
    \item \textbf{Cross-Brand SKU Collisions:} The retrieval system occasionally matched highly similar product specifications across different brands (e.g., matching a generic ``30W LED Panel'' to the wrong manufacturer) when the source document omitted explicit brand identifiers.
    \item \textbf{LLM Schema Deviation:} In approximately 1.5\% of cases, the generative model deviated slightly from the strict JSON schema (e.g., hallucinating nested objects for bundled promotions), causing downstream parsing exceptions before reaching the resolution module.
\end{itemize}